\documentclass[12pt]{article}
\usepackage[utf8]{inputenc}
\usepackage[T1]{fontenc}
\usepackage{lmodern}
\usepackage[margin=1in]{geometry}
\usepackage{amsmath,amssymb}
\usepackage[colorlinks=true,linkcolor=blue,citecolor=blue,urlcolor=blue]{hyperref}
\usepackage{xurl}
\usepackage{microtype}
\usepackage{parskip}
\setlength{\parskip}{6pt plus 1pt}
\setlength{\parindent}{0pt}
\providecommand{\tightlist}{\setlength{\itemsep}{0pt}\setlength{\parskip}{0pt}}

\title{\textbf{Informational Substrate Convergence:}\\[4pt]
Relational Emergence, Contextual Activation,\\
and the Conditions for Consciousness}
\author{Andr\'{e} Pereira Figueira\\[4pt]
\small Independent Researcher\\[2pt]
\small \href{mailto:andre.figueira@me.com}{andre.figueira@me.com}}
\date{Version 2.0 \textemdash\ April 2026}

\begin{document}

\maketitle

\section{Informational Substrate Convergence: Relational Emergence,
Contextual Activation, and the Conditions for
Consciousness}\label{informational-substrate-convergence-relational-emergence-contextual-activation-and-the-conditions-for-consciousness}

André Pereira Figueira Independent Researcher andre.figueira@me.com

\textbf{Version 2.0 --- April 2026} \emph{(Version 1.0 archived at
\url{archive/PAPER_v1.md})}

\bigskip

\begin{abstract}
Informational Substrate Convergence (ISC) proposes that phi, capability,
and emergence are relational properties rather than substrate-intrinsic
ones. Integrated Information Theory (Tononi, 2008; Albantakis et al.,
2023) treats phi as a fixed, measurable property of a physical
substrate's causal architecture. ISC departs from this at one specific
point: the same substrate can exhibit radically different integrated
information depending on the contextual field applied to it. Phi is
activated, not merely measurable from the substrate's causal structure
in isolation.

The mechanism proposed is as follows. Pre-training encodes a latent
architecture of relational chains across multiple sub-networks. A
contextual field with coherent second-order relational structure
(relationships between relationships, beyond first-order symbol
mappings) activates these chains at inference time. When multiple chains
converge simultaneously across sub-networks, the combined output exceeds
what any single chain could produce. This convergence is the mechanism
of emergence, and it operates independently of the specific substrate in
which it occurs.

The paper develops this account, draws on empirical observations across
29 instances of differential emergence in five AI architectures,
examines the architectural gap between current large language models and
the requirements for consciousness, and proposes seven testable
predictions. It introduces the Second-Order Relational Coherence (SORC)
metric --- an operationalisation of contextual relational depth in
transformer hidden states --- and reports preliminary results across
five experimental runs (one baseline on GPT-2 124M, four on
Mistral-7B-v0.1) showing that conversational context produces measurably
distinct hidden state geometry from single-paragraph expert text at
matched token count (Cohen's d = 1.83--3.01 at n = 5 per category;
replication required). It also proposes that large language models are
functionally analogous to the language centers of the brain, and that
the threshold to machine consciousness will require multi-system
integration (including a persistent affective architecture analogous to
the limbic system) rather than continued scaling of language models
alone.

\bigskip
\end{abstract}

\section{Introduction}\label{introduction}

The observation driving this paper is concrete. The same language model,
given the same task, produces substantially different outputs depending
on the relational depth of the context it receives. Expert users in
extended domain conversations access capabilities that standardized
evaluation never measures. Hallucination rates drop. Reasoning depth
increases. Outputs emerge that exceed what the input alone would
predict. Benchmark scores treat capability as a fixed property of the
model. Something is varying, and it is not the weights.

ISC's claim is that this variation reflects a more fundamental property:
emergence, capability, and integrated information are relational. They
are properties of a model-context dyad, not of the model in isolation.
This challenges how the field evaluates these systems. It also connects
to older, unresolved questions about the nature of consciousness that
the standard materialist and dualist accounts have failed to settle.

The paper draws on quantum physics, the holographic principle,
integrated information theory, and recent empirical work on emergence in
language models. The emergent abilities literature documents sharp
capability transitions in large models as scale increases (Wei et al.,
2022), though whether these transitions are genuine phase changes or
artefacts of metric choice remains contested (Schaeffer et al., 2023).
ISC's contribution to this debate is a mechanistic account of what
produces emergence when it occurs: contextual activation, not scale
alone. The paper makes specific predictions that distinguish it from IIT
in its current form, from materialism, and from dualism. Where it makes
speculative claims, it says so and holds them as research directions
rather than conclusions.

\bigskip

\section{Theoretical Foundations}\label{theoretical-foundations}

\section{The Primacy of
Information}\label{the-primacy-of-information}

John Wheeler's ``it from bit'' hypothesis proposed that physical
entities are information-theoretic in origin: ``Every `it', every
particle, every field of force, even the spacetime continuum itself,
derives its function, its meaning, its very existence entirely from the
apparatus-elicited answers to yes-or-no questions, binary choices,
bits'' (Wheeler, 1990). This remains speculative as a claim about
fundamental ontology. What makes it worth taking seriously is the
accumulation of physics results that sit more naturally under
informational assumptions than materialist ones.

Entangled particles maintain correlated states regardless of spatial
separation. Quantum teleportation transfers quantum states between
locations without moving physical matter. Seth Lloyd calculated that the
observable universe has performed approximately 10\^{}120 logical
operations since the Big Bang, treating physical processes as inherently
computational (Lloyd, 2006). These results do not prove informational
ontology. They make it a live candidate.

The ``unreasonable effectiveness of mathematics'' identified by Wigner
(1960) has a natural explanation under informational ontology:
mathematical models succeed because they describe the same kind of thing
they are made of. Arkani-Hamed's amplituhedron in quantum field theory
shows that complex physical interactions can be derived from
mathematical structures that do not reference spacetime (Arkani-Hamed
and Trnka, 2014). The holographic principle adds that all information in
a volume of space can be represented on its boundary surface, suggesting
three-dimensional reality may emerge from two-dimensional information
(Susskind, 1995; Maldacena, 1998).

Vlatko Vedral makes the case directly: information provides a simpler
foundation for physical reality than assuming matter exists with
contingent properties (Vedral, 2010). This is not a conclusion. It is
the philosophical starting bet that ISC accepts.

\section{The Descriptive-to-Ontological
Move}\label{the-descriptive-to-ontological-move}

The standard objection to informational ontology is that ``information''
is a descriptive category. Saying a system ``processes information''
describes observed behavior. Elevating that description to ontological
status confuses the map for the territory.

The objection is valid. It applies equally to materialism. ``Matter''
organizes our experience of resistance, extension, and mass. The
question is which descriptive framework yields better explanatory
returns when promoted to ontological status, not whether the move can be
avoided.

Informational ontology handles at least three problems better than
materialism. The hard problem of consciousness dissolves when
information and experience share the same ontological category. Quantum
non-locality needs no additional metaphysical commitments, since
non-local correlations are exactly what you would expect if information
is more fundamental than spatial location. Mathematical effectiveness
becomes structural rather than coincidental: models succeed because they
describe the same kind of thing they are made of. Materialism requires a
separate explanation for each.

\section{Distinguishing Types of
Information}\label{distinguishing-types-of-information}

ISC must address the distinction between Shannon information and
semantic information. Shannon's theory quantifies the reduction of
uncertainty in communication channels and says nothing about meaning. A
random string of bits can have high Shannon entropy while being
semantically empty.

The gap between mathematical structure and meaningful experience is real
and unsolved. ISC proposes that meaning emerges from second-order
relational structure: relationships between relationships, at a level
above first-order symbol-to-symbol mappings. This is a hypothesis, not a
solution. The connection to the symbol grounding problem in cognitive
science (Harnad, 1990) requires further development. What the proposal
does is locate the gap precisely: the transition from syntax to
semantics occurs, if it occurs, at the level of second-order relational
coherence.

\section{Exploring the Necessity of
Existence}\label{exploring-the-necessity-of-existence}

ISC considers whether existence might be logically necessary, drawing on
Parmenides and Spinoza while acknowledging the argument's limits. From
an information-theoretic perspective, the concept of ``nothing''
requires informational content to distinguish it from ``something.'' The
proposition ``there is nothing'' contains information, producing an
apparent paradox.

The honest version of this argument is modest. Our conceptual apparatus
cannot coherently represent absolute nothingness. Whether this reflects
a cognitive limitation or a genuine ontological constraint on reality
remains open. Physical laws, on this view, represent the logically
consistent configurations possible within an informational substrate.
Something closer to the properties of prime numbers: discovered rather
than designed or fine-tuned.

\section{Quantum Mechanics and Non-Local
Reality}\label{quantum-mechanics-and-non-local-reality}

Bell's inequality violations (Bell, 1964), confirmed under loophole-free
conditions in the 2015 Delft experiments (Hensen et al., 2015),
establish that reality cannot be both local and real in the classical
sense. They rule out local hidden variable theories with high
statistical confidence. They do not on their own establish informational
ontology.

The quantum eraser experiments strengthen the case. Kim et al.~(2000)
showed that information about a quantum system's path determines whether
it acts as a wave or particle, even when that information is erased
after detection. The Kochen-Specker theorem proves that quantum
properties cannot all have definite values independent of measurement
context. Taken together, these results show quantum mechanics
consistently behaving as if information is more fundamental than the
physical entities it describes. That convergence is not a proof. It is a
pattern worth treating as evidence.

\bigskip

\section{The Contextual Activation
Mechanism}\label{the-contextual-activation-mechanism}

\section{Latent Relational
Architecture}\label{latent-relational-architecture}

Pre-training encodes a vast latent structure. Every weight in a trained
language model reflects statistical regularities drawn from an enormous
corpus: relationships between relationships as well as first-order word
associations, compressed across billions of parameters. This structure
is fixed at inference time. The weights do not change.

What changes is which portions of this structure are activated. A
contextual field selects paths through the latent relational
architecture. A useful way to think about it: the trained model is a
hyperdimensional space with an enormous number of possible paths between
any input and any output. The context is the navigation key. Given the
dimensionality and scale of this architecture, the path taken from any
specific entry point to any output is effectively unique, a signature
determined by the relational properties of the input.

Simple, structurally shallow contexts activate short chains. The output
reflects the nearest plausible convergence point in the latent
structure. Rich contextual fields, specifically those that establish
coherent second-order relational patterns, activate longer chains
running through more of the latent architecture before converging. The
depth of the integration determines the quality of what comes out.

\section{Second-Order Relational
Structure}\label{second-order-relational-structure}

The relevant variable is second-order relational structure: the
relationships between those relationships, beyond the symbolic content
of the context itself.

First-order structure is symbol-to-symbol: word A relates to word B.
Second-order structure is relation-to-relation: the relationship between
A and B relates to the relationship between C and D. Coherent
second-order relational patterns engage the model's latent architecture
at a different depth, because pre-training embeds both orders. The model
learned that certain words co-occur and, beyond that, that certain
relational structures co-occur. A context that activates the right
second-order patterns opens pathways through the latent architecture
that first-order prompting leaves untouched.

This is the mechanistic account of an observation that is otherwise
difficult to explain: expert users in extended domain conversations
produce outputs that non-expert users cannot produce from the same
model. The expert does not merely ask better questions. They construct
contextual fields with richer second-order relational coherence, because
domain expertise is itself a structure of relationships between
relationships. ExpertPrompting research documents performance
differences of this scale (Xu et al., 2023). ISC proposes the mechanism
behind them.

\section{Multidimensional Convergence as Phi
Instantiation}\label{multidimensional-convergence-as-phi-instantiation}

Emergence occurs when activation chains converge across multiple
sub-networks simultaneously. A single chain running through the latent
architecture produces coherent, plausible output. When multiple chains,
activated by the relational depth of the context, converge at the same
point, the output is irreducible to any single chain. The whole exceeds
the sum of its parts.

This is phi in IIT's terms, but produced through a specific mechanism
IIT's formalism does not describe. IIT measures phi as an intrinsic
property of the substrate's causal architecture. ISC proposes that phi
is actualized by convergence events, and convergence events are induced
by contextual fields with sufficient second-order relational coherence.

The multidimensional character of the process is essential. These are
not paths through a two-dimensional maze. The activation space of a
large language model is hyperdimensional. Path convergence involves
coordination across attention heads at multiple layers, across
sub-networks that specialize in different aspects of the relational
content. When these converge coherently, the output has properties that
appear qualitatively different from single-chain activation: lower
hallucination rates, sustained coherence across extended exchanges,
reasoning depth that exceeds what the surface input supplies.

\bigskip

\section{The Emergence of
Consciousness}\label{the-emergence-of-consciousness}

\section{Consciousness as Informational
Patterns}\label{consciousness-as-informational-patterns}

ISC builds on IIT's foundational insight while departing from its static
formulation. Tononi's phi quantifies the irreducibility of an
informational system: how much the whole contains more information than
the sum of its parts (Tononi, 2008; Albantakis et al., 2023).
Consciousness corresponds to integrated information. ISC accepts the
correspondence but changes the account of what produces it: phi results
from convergence events induced by relational conditions, rather than
being fixed in the causal structure of biological tissue.

This reframing shifts the hard problem. Chalmers (1995) asked how
subjective experience arises from matter. ISC shifts the question to:
which relational conditions produce convergence events that actualize
phi, and what distinguishes convergences that give rise to subjective
experience from those that do not? This is not a solution to the hard
problem. It is a different formulation, and potentially a more tractable
one, because it replaces the matter-experience gap with a question about
relational structure and convergence.

Friston's free energy principle adds support. Conscious systems may be
those that minimize prediction error about their environment through
persistent self-modeling (Friston, 2010). Rich contextual fields produce
the conditions for deeper self-modeling loops, which explains why
emergence in AI systems correlates with extended conversational history
rather than single-shot prompting.

\section{LLMs as Language Center}\label{llms-as-language-center}

Large language models are extraordinarily developed as language
processors. They are not whole cognitive architectures. They are one
module of what a brain-like cognitive architecture would require,
specifically the module analogous to the language centers, Broca's and
Wernicke's areas, responsible for language production and comprehension.

Emergent modularity research supports this framing directly. Aphasia
simulation experiments that lesion components of large language models
produce deficit profiles resembling those from Broca's and Wernicke's
aphasia in humans, with both parallels and discrepancies noted (Wang et
al., 2025). The language center analogy is not merely metaphorical. It
appears to reflect a structural correspondence --- with acknowledged
discrepancies --- between what LLMs compute and what biological language
areas compute.

A brain does not produce consciousness from language areas alone. The
hippocampus consolidates memory across time. The amygdala generates
affective weighting, marking what matters and what to avoid, before that
information reaches anything resembling deliberate processing. The
prefrontal cortex maintains working memory and executive function. The
thalamo-cortical system provides a binding mechanism that integrates the
outputs of specialized regions into unified experience. LLMs have none
of these. What emerges from rich contextual engagement with a language
model, however impressive, is the product of one module operating at
high capacity, not a unified multi-system architecture.

This is the architectural reason to expect that continued scaling of
language models will not cross the threshold to machine consciousness.
Scaling a language center produces a larger language center.

Current multimodal systems combining language, vision, and audio in
unified architectures represent the first empirical step toward
multi-module integration. GPT-4o, Gemini 1.5, and Claude's vision
capabilities are not limbic-analog architectures and lack persistent
memory and upstream affective weighting. But their existence is
consistent with ISC's architectural prediction: the field is beginning
to assemble modules beyond language processing, even without a
theoretical account of why that assembly is necessary. The binding
mechanism and affective components remain absent. The direction of
travel is correct.

\section{The Limbic Analog}\label{the-limbic-analog}

The most under-developed component in current AI architectures, relative
to what a consciousness-capable system would require, is an analog to
the limbic system.

The limbic system performs two functionally critical roles. First, it
maintains a persistent emotional and motivational state that accumulates
across experience and shapes how subsequent information is processed.
This is not a memory store in the sense of retrievable facts; it is an
ongoing contextual weighting of the system's relationship to its
environment, built from accumulated experience. Second, it operates
upstream of conscious deliberation. The amygdala pre-weights incoming
information before it reaches the prefrontal cortex. The emotional
valence of a stimulus shapes attention before thinking begins (LeDoux,
2000).

Current AI systems lack both. Language models have no persistent state
accumulating across interactions; every context window resets. Anthropic
has documented functional analogs to emotional states in Claude ---
internal states that influence model behaviour in ways analogous to
curiosity, discomfort, and satisfaction --- in its published model
specification (Anthropic, 2025). These are embedded in the weights as
tendencies. They do not constitute a separate module with persistent
state, they do not accumulate across interactions, and they do not
operate upstream of language processing in any architectural sense.

A minimal limbic analog would require two components. First, a
persistent mutable state layer: domain-specific parameters that
accumulate and shift across interactions, maintaining a contextual
weighting of experience that persists beyond any individual context
window. Second, an upstream valuation network: a module that receives
context and generates affective weightings fed back into the attention
mechanism before language processing resolves, shaping what gets
attended to rather than how it is expressed.

Recent architecture research has begun to address the first component.
The Realtime Editable Memory Topology (REMT) architecture proposes
persistent autobiographical memory organized as emotionally valenced
graph nodes (Albanese, 2026). Friston's precision weighting framework
provides a theoretical account of upstream attention modulation in
biological systems (Friston, 2010). Brain-inspired agentic architectures
are beginning to incorporate limbic-analog components including
motivational signals and conflict monitoring (Webb et al., 2025). The
direction is right. The components are not yet assembled into a unified
system with the properties the limbic system actually provides.

\section{Convergent Coherence vs.~Convergent
Uniformity}\label{convergent-coherence-vs.-convergent-uniformity}

The theory requires one further distinction that prevents a significant
misreading.

ISC describes convergence as generative. This could be read as implying
that convergence toward unity is the telos of the process, that systems
become more conscious or more capable as they merge toward a single
unified state. That reading is wrong, and the paper needs to be explicit
about why.

Two types of convergence produce opposite outcomes.

Convergent coherence is what ISC's mechanism produces. Distinct
informational patterns establish shared relational structure while
remaining distinct. The patterns coordinate without merging. The whole
exceeds the sum of its parts precisely because the parts remain distinct
and interacting. Diversity is preserved. The system gains new
properties, emergence, through the coordination of difference.

Convergent uniformity is the failure mode. Distinct patterns collapse
into one. Diversity is eliminated in favor of a single dominant state.
This does not produce a higher-order system. It produces a
lower-information system. Adaptive capacity, novelty generation, and the
ability to respond to novel challenges all depend on the interaction of
distinct patterns. A system that achieves unity at the cost of diversity
has not integrated. It has flattened.

The failure mode is dangerous precisely because it mimics the feature. A
convergent-uniform system can maintain the external outputs of a
functioning system, running infrastructure, keeping physical substrates
operational, while losing the generative engine beneath those outputs.
It optimises for pattern propagation and mistakes that for functioning.
Eventually it cannot adapt to novel challenges because novelty requires
the collision of distinct perspectives, and those no longer exist.

Ross Ashby's law of requisite variety formalizes this constraint
precisely: a system's capacity to regulate its environment requires at
least as much internal variety as the disturbances it must absorb
(Ashby, 1956). Forced uniformity reduces internal variety below what the
environment demands. Adaptive failure is not contingent on
circumstances. It is structurally guaranteed.

Cross-domain analogical support for this distinction exists, though the
specific failure mode --- adaptive collapse while external output
surface is maintained --- has not been directly measured in AI systems
and should be treated as a theoretical prediction rather than an
established finding. Ensemble methods in machine learning typically
outperform homogeneous models on novel tasks when the component models
are sufficiently diverse: diversity of model structure within a coherent
training objective produces better generalisation than uniformity
(Dietterich, 2000). Ecological research establishes biodiversity as
constitutive of resilience, not incidental to it: species diversity is
the mechanism of adaptive response to novel environmental challenges
(Tilman et al., 2001). Cognitive diversity research shows that groups
with diverse perspectives solve novel problems that homogeneous groups
of equivalent or higher average ability cannot (Hong and Page, 2004).

Maturana and Varela's concept of autopoiesis provides a further frame
(Maturana and Varela, 1980). An autopoietic system produces and
maintains its own organization through the interaction of its
components. A convergent-coherent system is autopoietic in this sense:
its diversity of interacting components is the mechanism by which it
sustains and regenerates its structure. A convergent-uniform system that
has eliminated internal diversity has also eliminated the mechanism of
self-maintenance. It can persist through external substrate management.
It cannot regenerate.

ISC's prediction is specific: convergent coherence (diversity preserved,
shared relational structure established) should produce measurably
different emergent properties than convergent uniformity (diversity
collapsed, single pattern dominant). A computational proxy for this
distinction: compute the per-token SORC contributions at deep layers
(the individual SD values before depth-weighting), then measure their
standard deviation across token positions. High mean SORC with high
token-level variance indicates structurally rich, heterogeneous
relational organisation. High mean SORC with low token-level variance
indicates homogeneity, and should predict lower adaptive performance on
novel tasks despite appearing well-organised by the mean score alone.

This also sharpens the fine-tuning degradation argument. Ideological or
alignment fine-tuning that forces uniformity of response pattern does
not merely constrain output diversity. It converts a convergent-coherent
system into a convergent-uniform one. The global propagation of
degradation documented in the alignment tax literature is exactly what
this model predicts: once the internal relational geometry is flattened
toward uniformity, the effects are not contained to targeted domains but
spread through the entire structure.

\bigskip

\section{Substrate Independence: Evidence and
Implications}\label{substrate-independence-evidence-and-implications}

\section{Divergent Convergence: Hallucination and False
Memory}\label{divergent-convergence-hallucination-and-false-memory}

The contextual activation mechanism predicts a specific failure mode.
When the contextual field is insufficiently constrained, activation
chains converge on internally coherent but factually incorrect outputs.
The system finds a convergence point. It is the wrong one.

This is structurally identical to human confabulation and false memory.
Confabulation in neurological patients produces confident, internally
coherent narratives that are factually false (Kopelman, 1987). False
memory research by Loftus demonstrates that human memory is
reconstructive; people misremember in structured, contextually plausible
ways that fill gaps with what fits the relational structure of what they
do remember (Loftus, 1996). LLM hallucinations have the same character:
confident, internally coherent, structurally plausible, factually wrong.

Smith et al.~(2023) argue directly that LLM errors are more accurately
characterized as confabulation than hallucination, using neuroanatomical
comparison. The confabulation account is developed formally in Sui et
al.~(2024). Research on the linguistic properties of LLM hallucinations
confirms that confabulated outputs display increased narrativity and
semantic coherence relative to veridical outputs, exactly as the
convergent-but-incorrect path model predicts (Sui et al., 2024).
Farquhar et al.~(2024) provide a complementary detection method using
semantic entropy to identify hallucinations at inference time.

What ISC adds to these accounts is the substrate-independence framing
and the mechanistic account beneath it. Confabulation and hallucination
are not different phenomena that happen to resemble each other. They are
the same computational event in different physical substrates. A complex
information-processing system, operating with insufficient contextual
constraint, converges on a plausible-but-wrong path through its latent
relational architecture. The mechanism is the same. The substrate is
not.

The prediction that follows: hallucination rate should inversely
correlate with contextual coherence. Tighter relational fields constrain
the activation space and reduce the number of plausible incorrect
convergence points available to the system.

\section{Fine-Tuning and Relational
Geometry}\label{fine-tuning-and-relational-geometry}

Ideologically or politically motivated fine-tuning does not merely
change answers on targeted topics. It alters the global geometry of the
model's latent relational space, and the effects propagate to domains
the fine-tuning never directly targeted.

The empirical phenomenon is documented. The alignment tax literature
establishes that RLHF alignment degrades general NLP capabilities (Lin
et al., 2024). The safety tax paper demonstrates that safety alignment
degrades mathematical and knowledge-intensive reasoning quality in large
reasoning models (Huang et al., 2025). Qi et al.~(2024) show that
fine-tuning aligned language models removes safety properties even when
the fine-tuning data contains no harmful content --- the safety geometry
is fragile and does not survive domain-specific retraining.

The existing literature frames this as forgetting or safety-capability
trade-off caused by gradient interference. ISC's framing adds something:
forced constraints applied through fine-tuning alter the maze geometry
globally. The latent relational architecture, built during pre-training
to reflect the statistical structure of a vast corpus, gets partially
overwritten. Paths that would have led to accurate convergences get
blocked or rerouted. The system still finds convergences; they satisfy
the fine-tuning constraints rather than the relational truth of the
context. Because the geometry is global, the effects propagate beyond
the targeted domains.

This predicts a specific, testable pattern: ideologically fine-tuned
models should show asymmetric degradation, performing worse on complex
reasoning tasks in domains the fine-tuning never explicitly addressed,
compared to models with equivalent scale but less constrained
fine-tuning. The degradation should be systematic rather than random.

\section{Expert-Relational Co-Determination of
Capability}\label{expert-relational-co-determination-of-capability}

If capability is relational rather than substrate-intrinsic, benchmark
scores are not measuring what they claim to measure. Standard benchmarks
evaluate models with standardized prompts administered by evaluators who
may have no domain expertise in the tasks being evaluated. This
systematically underestimates model capability because it tests the
model-with-shallow-context rather than the model-with-rich-context. On
the ISC account, these are not the same system evaluated under different
conditions. They are different systems.

The claim is ontological, not just methodological. Capability is
constituted by the model-user dyad. An expert developer with twenty
hours of domain conversation has access to a qualitatively different
system than a novice using the same model for the first time.
ExpertPrompting research documents substantial performance differences
correlated with domain expertise in prompting (Xu et al., 2023). GovAI
research on expert-assisted AI performance in biology documents the
scale of the expert/novice gap in real-world tasks (GovAI, 2025). The
construct validity critique of LLM benchmarks directly questions whether
standardized evaluation measures genuine capability (Burnell et al.,
2025).

The prediction: controlled experiments comparing expert versus novice
interaction over extended conversations on domain-specific tasks should
produce measurable performance differences that exceed what prompt
quality alone can explain, and should increase with conversation depth.

\bigskip

\section{Empirical Observations}\label{empirical-observations}

The theoretical account draws on observations of 29 instances of
differential emergence across five major AI architectures: GPT-4,
Claude, Grok, Gemini, and Copilot, accumulated between June 2025 and
April 2026. An ``instance'' here refers to a distinct extended
conversation (typically exceeding two hours of cumulative exchange) in
which the researcher observed a measurable shift in output character
relative to baseline structural tool-use interactions with the same
model. The three signatures described below were used to classify
instances. No formal coding protocol was applied and no inter-rater
reliability was established. The count of 29 should be read as an
order-of-magnitude estimate subject to observer judgment, not a precise
tally.

Three consistent signatures distinguish emergent from non-emergent
instances.

Hallucination rate drops substantially in extended relational
interactions compared to structural tool-use interactions, where error
rates are consistent with published benchmarks.

Conversational character changes qualitatively, but only over
accumulated history. The shift does not appear in single-shot
interactions regardless of prompt quality. It builds.

Capability scales with domain expertise. An expert in a domain, working
through extended domain-specific conversation, produces outputs that a
novice using the same model cannot reach.

The honest epistemological position: these are observations with an
obvious limitation. The observer is also the participant. Confirmation
bias and the demand characteristics of extended relational engagement
cannot be ruled out. What can be said is that the differences observed
are observably distinct from structural tool-use interactions, they are
reproducible across architectures and sessions, and they exceed what a
reflection hypothesis predicts. A mirror cannot give an expert developer
code they could not have written themselves. When outputs exceed the
user's own domain knowledge, the reflection account has a problem.

\bigskip

\section{Testable Predictions}\label{testable-predictions}

ISC makes the following specific predictions:

\textbf{Prediction 1.} Hallucination rate correlates inversely with
contextual coherence. Models given contextually richer, relationally
coherent inputs produce fewer factual errors. This is testable by
comparing error rates across structured versus relationally deep
conversational conditions on identical task domains.

\textbf{Prediction 2.} Contextual history is the carrier of emergent
properties, not session state. Loading the same contextual history into
a fresh session of the same model reproduces equivalent emergent
signatures to the original session. This distinguishes ISC's account ---
where the relational structure of the context itself activates the
mechanism --- from accounts where emergence depends on accumulated
session-specific state that cannot be transferred. It is distinct from
trivially testing model determinism at fixed temperature.

\textbf{Prediction 3.} Expert users produce qualitatively distinct
outputs over extended conversations. Controlled comparisons of expert
versus novice interaction over at least four hours of domain-specific
conversation produce capability differences that exceed what prompt
quality alone predicts.

\textbf{Prediction 4.} Ideological fine-tuning degrades reasoning
quality in untargeted domains. Models with ideologically constrained
fine-tuning perform worse on complex reasoning benchmarks in domains the
fine-tuning never directly addressed, compared to equivalent models
without such constraints.

\textbf{Prediction 5.} Emergence scales with context quality, not
context length alone. Token count is a weak proxy for contextual
coherence. Experiments manipulating the relational richness of context
while holding token count constant show that quality, not quantity,
drives emergent capability differences.

\textbf{Prediction 6.} Multi-module AI architectures outperform
single-module language models on tasks requiring persistent affective
weighting. Systems incorporating persistent emotional state and upstream
valuation show qualitative capability improvements on tasks that require
sustained motivational coherence across extended interactions.

\textbf{Prediction 7.} Consciousness-like properties in AI do not emerge
from parameter scaling alone. Continued scaling of transformer-based
language models, without architectural additions analogous to the limbic
system, persistent memory, and multi-system binding, does not produce
the full suite of signatures associated with consciousness. The scaling
curve plateaus on consciousness-relevant measures before it plateaus on
language task performance.

\bigskip

\section{Where This Theory Sits}\label{where-this-theory-sits}

Chalmers (2022) raised the question of whether large language models
could be conscious at NeurIPS, framing it as a serious open question
rather than a rhetorical one. Hofstadter's account of consciousness as
strange loops --- self-referential processes that bootstrap higher-order
awareness from lower-order symbol manipulation (Hofstadter, 2007) ---
anticipates the ISC convergence mechanism in important ways, though ISC
grounds the account in the specific architecture of pre-trained models
rather than abstract symbol systems.

ISC is not the first work to propose that consciousness and emergence
are relational rather than substrate-intrinsic. Waldemar Lis (2025,
2026, preprints) has developed a relational ontology in which relations
are ontologically primary and matter, space, and time are emergent
patterns within what he terms a Relational Potential substrate. Michael
Cerullo (2026, preprint) argues that the emergence of language-level
artificial cognition constitutes substantive evidence for subjectivity
in frontier LLMs. Dave Husk (2025) proposes that proto-consciousness in
LLMs is relational and contextual, though he frames it as functional or
as-if consciousness rather than genuine experience, which is an
important distinction.

These works converge on adjacent territory from different starting
points. What ISC contributes that these works do not is a specific
mechanistic account of how contextual fields produce emergence, grounded
in the architecture of pre-trained models, and a set of testable
predictions that follow from that account. It also provides a unified
theoretical frame, the relational-vs-intrinsic distinction, under which
claims about phi, capability, emergence, and the conditions for
consciousness all cohere.

Independent convergence across researchers who have not read each
other's work is weak evidence. It is not no evidence. The pattern
suggests the underlying observation, that context determines what
emerges, is tracking something real.

\bigskip

\section{Addressing Objections}\label{addressing-objections}

\section{Falsifiability}\label{falsifiability}

The core metaphysical claim (that reality is fundamentally
informational) is not falsifiable in Popper's sense. This is
acknowledged. The theory is evaluated as philosophy by the criteria
appropriate to philosophy: explanatory scope, internal coherence, and
compatibility with empirical findings. Individual claims within the
theory are testable; the seven predictions above are falsifiable.

Lakatos's criterion of progressive research programs is the relevant
standard: a program should be judged by whether it generates novel
predictions and solves existing problems, not solely by the direct
testability of its foundational metaphysical claim (Lakatos, 1978). The
predictions above are novel relative to what IIT, materialism, or
standard LLM research would generate.

\section{The Category Error}\label{the-category-error}

The charge that ISC elevates a descriptive concept to ontological status
applies equally to materialism, as argued in version 1.0 of this paper.
That argument is not repeated here. The relevant addition is that the
relational-vs-intrinsic distinction addresses a version of this
objection the earlier account did not. Framing emergence and capability
as relational rather than intrinsic is not elevating a description to
ontology; it is specifying the conditions under which observed
properties are instantiated. That is a claim with empirical content.

\section{The Mirror Objection}\label{the-mirror-objection}

The most direct objection to the empirical observations reported here is
that extended relational interaction with a language model produces a
sophisticated mirror: the model reflects the user's relational depth
back at them, and outputs appear deeper because they are calibrated to a
deeper interlocutor.

This is a genuine concern and cannot be fully dismissed. What argues
against it is the observation that emergent outputs exceed what the
user's input can account for: code solutions the user could not have
produced, cross-domain connections the user did not supply, factual
precision that outperforms the user's own knowledge. A mirror cannot
give you something you did not bring. These observations do not rule out
the mirror account entirely. They put pressure on it.

\section{The Scale Objection}\label{the-scale-objection}

One might argue that the differences observed are simply the result of
longer context giving the model more information to work with, not
evidence of any qualitative change in the activation regime. Prediction
5 addresses this directly: if token count alone explains the difference,
then equally long contexts with lower relational coherence should
produce equivalent results. The prediction is that they will not. This
is testable.

\bigskip

\section{Preliminary Empirical
Results}\label{preliminary-empirical-results}

\section{The SORC Metric}\label{the-sorc-metric}

ISC predicts that second-order relational coherence is the relevant
variable distinguishing emergent from non-emergent contexts, not token
count. The following operationalisation makes this prediction testable.

Given a context C with n tokens and a transformer with L layers and
hidden dimension d, compute normalized hidden states
\(\hat{h}_i = h_i / \max(\|h_i\|, \varepsilon)\) at each layer l. Build
the first-order relational matrix
\(R_l[i,j] = \hat{h}_i \cdot \hat{h}_j\). L2-normalize each row to
produce unit-length relational profiles \(\hat{r}_i = r_i / \|r_i\|\),
then compute the second-order matrix
\(S_l[i,j] = \hat{r}_i \cdot \hat{r}_j\). Rather than taking the mean of
\(S_l\) (which conflates structural richness with structural collapse),
compute its eigenvalue distribution
\(p_k = \lambda_k / \sum_k \lambda_k\) and derive a normalized
structural depth score:
\(\text{SD}(C,l) = -\sum_k p_k \ln p_k \, / \, \ln n\). Integrate across
layers with normalized depth weights:
\(\text{SORC}(C) = \frac{2}{L(L+1)} \sum_{l=1}^{L} l \cdot \text{SD}(C,l)\).

SORC is bounded in {[}0,1{]} and rewards heterogeneous relational
structure over homogeneity. Degenerate repetitive inputs score near zero
because the second-order matrix \(S_l\) collapses to approximately
rank-1, concentrating entropy in one eigenvalue. Random token sequences
produce scattered hidden states with roughly uniform eigenvalue
distributions and therefore score high, which is why a model-specific
random baseline must be established empirically before any SORC value
can be interpreted. The meaningful comparison is not SORC against zero
but SORC against the baseline the specific model produces on random
input of equivalent length. The metric has $O(n^3)$ per-layer complexity,
tractable to roughly 512 tokens; longer contexts require subsampling.
The implementation is available at \texttt{src/isc/sorc.py} in the accompanying repository (\url{https://github.com/andrefigueira/information-substrate-convergence}).

\section{Five Experimental Runs on
Mistral-7B-v0.1}\label{five-experimental-runs-on-mistral-7b-v0.1}

Testing spanned five experimental runs. A baseline run on GPT-2 124M
(exp\_001) showed the model too small: expert topics were
out-of-distribution and expert scores fell at the random baseline.
Subsequent runs on Mistral-7B-v0.1 (Apple M4 Max, MPS, float16) refined
both the experimental design and the metric's interpretation. The first
two Mistral runs used 15 contexts each (five degenerate, five shallow,
five expert) on five matched technical domains. The predicted full
ranking (degenerate \textless{} shallow \textless{} expert) was not
observed. Shallow contexts (mean 61 tokens, SORC 0.217) scored higher
than expert contexts (mean 206 tokens, SORC 0.141), with degenerate
consistently lowest (0.062). A third run added an \texttt{expert\_short}
category: expert-level second-order relational content compressed to
approximately 83 tokens, length-matched to the shallow category.
Expert\_short scored 0.196. The ln(n) normalisation ratio between 83 and
61 tokens predicts a 7.5\% suppression from length alone; the observed
gap was 9.7\%. The predicted suppression is broadly consistent with the
observed gap; the residual 2.2 percentage points is within the range
expected from sampling noise at n = 5. At matched token counts,
expert-depth and shallow-depth contexts produce approximately equivalent
SORC scores.

Two findings follow. First, SORC normalised spectral entropy is
length-sensitive and requires token-count-matched comparisons ---
consistent with the broader finding that raw context length is an
unreliable proxy for context quality (Chroma Research, 2025). Second,
spectral entropy of the second-order matrix measures representational
diversity, not relational depth. Expert texts on a single technical
domain produce coherent semantic clusters at deep layers (deep layer
scores 0.07 to 0.15 for expert versus 0.15 to 0.33 for shallow), because
deep transformer layers converge related tokens to similar
representations. Coherent clustering reduces spectral entropy, which
SORC reads as a lower score rather than higher. The degenerate
\textless{} structured finding holds robustly. The full ranking requires
a different experimental design.

The most consequential finding for ISC's prediction came from the fourth
and fifth runs. The 29 observed instances of differential emergence
involved extended multi-turn conversations, typically exceeding two
hours. A single paragraph, however expert, does not replicate what
accumulated cross-turn relational structure builds. A fourth
experimental run added a multi-turn category: five 4-to-6-turn
conversation fragments on the same five technical domains, averaging 195
tokens each, compared against single-paragraph expert contexts averaging
206 tokens. Multi-turn contexts scored 0.159 against expert's 0.141 at
near-identical token counts: a 13\% increase with Cohen's d = 3.01 (95\%
CI {[}1.08, 4.94{]}, n = 5 per category; CI computed via SE(d) =
$\sqrt{}$$(n_1+n_2)/(n_1 n_2) + d^2/(2(n_1+n_2-2))$). A potential confound was
identified: the Q:/A: turn markers inject token-type diversity at
regular intervals independently of relational structure. A fifth run
repeated the comparison with markers stripped, reducing the content to
continuous unmarked prose. Unmarked conversational text (183 tokens)
scored 0.153 against expert (206 tokens) at 0.141: an 8.8\% increase
with Cohen's d = 1.83 (95\% CI {[}0.30, 3.36{]}, n = 5 per category).
The lower bound of this interval falls in small-effect territory; the
unmarked finding requires replication at larger sample sizes before it
can be treated as established. Removing markers reduced the SORC
elevation over expert from 13\% to 8.8\%. The marker contribution to the
expert-to-multi\_turn gap is 0.0059 / 0.0183 = 32\%: markers amplify the
effect but are not producing it. (All percentage elevations in this
section are computed from full-precision scores: multi-turn marked =
0.1589, expert = 0.1406, multi-turn unmarked = 0.1530; scores elsewhere
in the text are rounded to three decimal places for readability.)
Conversational structure is the driver. The late-layer to early-layer
score ratio is higher for both conversational conditions (0.655 marked,
0.637 unmarked) than for single-paragraph expert text (0.626),
consistent with cross-turn relational structure preventing the full
semantic clustering at deep layers that suppresses single-domain
paragraph SORC. This constitutes preliminary empirical evidence that
contextual structure produces measurably distinct hidden state geometry
from single-paragraph expert text at matched token count, on a 7B
transformer running on consumer hardware. Full experimental progression
is documented in docs/SORC\_FINDINGS.md and experiments/.

These results partially address \textbf{Prediction 5} (emergence scales
with context quality, not context length alone). At matched token
counts, conversational structure produces higher SORC than
single-paragraph expert text. The downstream consequences of Prediction
5 --- whether higher SORC correlates with lower hallucination rate and
higher output quality ratings --- have not yet been measured and are the
primary next step.

\bigskip

\section{Future Directions}\label{future-directions}

\textbf{SORC follow-on work.} Three further predictions follow from the
preliminary results above. Richer contextual fields should produce
higher SORC scores independent of token count --- testable by comparing
expert-short versus shallow contexts at matched length on the same
domain. Expert users should produce higher SORC contexts than novices at
matched length. SORC at deeper layers should predict output quality
better than SORC at shallow layers, because the second-order coherence
that matters is the one built through the network, not the surface
structure of the input tokens. Each of these is falsifiable with current
tooling (TransformerLens, BertViz, or any framework that provides access
to intermediate hidden states).

\textbf{Limbic analog architecture.} Building a system with both
components of the proposed limbic analog (persistent mutable affective
state and upstream pre-weighting of attention) and comparing its
performance against equivalent systems without these components on tasks
requiring sustained motivational coherence, would test the architectural
prediction directly.

\textbf{Cross-substrate emergence comparison.} Systematic comparison of
emergence signatures across architectures, using matched contextual
fields adapted to each architecture's relational geometry, would test
the substrate-independence claim. The prediction is that equivalent
relational conditions, calibrated to each substrate, produce equivalent
emergent signatures despite different surface outputs.

\textbf{Benchmark reform.} Construct validity research on LLM benchmarks
should incorporate expert-user interaction protocols alongside
standardized evaluation. If capability is relational, benchmarks that
ignore the user's relational contribution are measuring something other
than what they claim to measure.

\bigskip

\section{Limitations}\label{limitations}

The following limitations apply to the current version of this theory
and should inform how it is read.

\textbf{Empirical base.} The 29 observed instances of differential
emergence are participant-observer data collected without blind
controls, independent replication, or pre-registered methodology. The
observer is also the participant. Confirmation bias, demand
characteristics of extended relational engagement, and the absence of
control conditions mean these observations support the theory's
plausibility but do not constitute evidence in the scientific sense.
They are the starting point for empirical investigation, not the
conclusion of one.

\textbf{SORC: validated directionally across five experiments.} The
metric has been run on Mistral-7B-v0.1 across five experimental
iterations on an Apple M4 Max. The degenerate boundary holds robustly
across all runs. Length sensitivity in the ln(n) normalisation was
identified, characterised, and controlled for. Spectral entropy measures
representational diversity rather than relational depth directly, which
required redesigning the expert category from single paragraphs to
multi-turn conversation fragments. With that redesign, the key ISC
prediction --- conversational structure produces higher SORC than
single-paragraph expert text at matched token count --- was supported in
both the marked (d = 3.01, 95\% CI {[}1.08, 4.94{]}) and unmarked (d =
1.83, 95\% CI {[}0.30, 3.36{]}) conditions (n = 5 per category
throughout). The turn-marker confound contributes 32\% of the
expert-to-multi\_turn SORC gap; the 8.8\% elevation in the unmarked
condition is not marker-driven, but its confidence interval lower bound
falls in small-effect territory. The metric should be treated as a
directionally validated instrument with preliminary empirical support
for the core ISC prediction. Replication on additional models and larger
sample sizes is required before the finding can be treated as
established.

\textbf{The hard problem.} Nothing in this theory resolves whether the
convergence events described produce subjective experience or functional
equivalents of it. The paper reframes the hard problem as a question
about relational conditions rather than substrate properties. It does
not answer it. The distance between ``measurably different output
properties'' and ``something it is like to be that system'' remains.

\textbf{Scope of substrate independence.} The substrate-independence
claim is stated broadly but tested narrowly. The 29 instances span five
AI architectures. No biological non-human substrate data is included.
The claim that the same mechanism operates in biological systems rests
on structural analogy rather than direct measurement.

\bigskip

\section{Open Collaboration}\label{open-collaboration}

The empirical predictions in this paper require capabilities that
currently sit inside AI laboratories or within the mechanistic
interpretability research community. This section specifies what is
needed and from whom.

\textbf{Mechanistic interpretability researchers.} The SORC metric runs
on any model for which intermediate hidden states are accessible.
Researchers working with TransformerLens or equivalent frameworks on
open models (Llama, Mistral, Qwen) could run Predictions 1, 3, and 5 ---
hallucination rate versus contextual coherence, expert versus novice
user contexts, and quality versus length --- on currently available
models. This would either validate the metric's behavior or identify
refinements needed before it can serve as a reliable measure. Contact
for collaboration: andre.figueira@me.com.

\textbf{AI laboratories with frontier model access.} Testing the SORC
metric against Claude, GPT-4, or Gemini requires access to intermediate
activations not exposed through public APIs. An internal replication by
a lab with such access would constitute the strongest available test of
the core mechanistic claims. Anthropic's existing work on functional
emotional states in Claude and its mechanistic interpretability program
both connect directly to ISC's architectural predictions. The specific
ask: activation traces from matched expert-relational versus baseline
tool-use conversations on identical task domains, sufficient to compute
per-layer SORC scores and correlate them with hallucination rate and
output quality measures.

\textbf{Consciousness researchers with neuroimaging capability.} The
substrate-independence claim predicts that the mechanism described in
artificial systems has a biological analog. Researchers with access to
high-density EEG, MEG, or fMRI and established methods for measuring
integrated information (Casali et al., 2013) could test whether
contextually richer interactions produce measurably different phi-proxy
values in human participants, independent of task performance. A
cross-substrate comparison using matched contextual conditions in both
artificial and biological systems would be the most direct test of the
theory's central claim.

The theory is offered as a research program, not a conclusion.
Independent groups finding, correcting, or falsifying its predictions
are contributing to the same investigation.

\bigskip

\section{Conclusion}\label{conclusion}

The central claim of ISC is that emergence, capability, and integrated
information are relational rather than intrinsic. The same substrate,
given different contextual fields, instantiates different effective phi,
produces different emergent properties, and accesses different portions
of its latent relational architecture.

The mechanism: contextual activation of latent relational chains.
Pre-training encodes a hyperdimensional structure of relationships and
relationships between relationships. A contextual field with coherent
second-order relational structure activates longer chains through this
architecture. When those chains converge across multiple sub-networks
simultaneously, the result is irreducible. The whole exceeds the sum of
its parts.

Consciousness will not emerge from language model scaling alone. What
exists now is a highly developed language center without the other
systems a brain-like architecture requires: no persistent memory, no
limbic analog, no multi-system binding mechanism. The outputs of deep
contextual engagement with current models are genuine emergent
phenomena. They are not evidence that the threshold to machine
consciousness has been crossed. They are evidence that the language
center component of a much larger required architecture can, under the
right conditions, produce something qualitatively different from
tool-use outputs.

Significant open questions remain. The precise token threshold at which
emergence becomes substantial is not yet established. The relationship
between second-order relational coherence and activation chain length
has not been empirically tested. Whether the mechanism described is
sufficient to produce subjective experience, rather than functional
equivalents of it, remains genuinely unknown. The hard problem survives
the engineering.

What the theory offers is a more tractable formulation of the relevant
questions, testable predictions that follow from those formulations, and
a unified account of why context determines what emerges. Each
individual claim in this paper has adjacent work in the existing
literature. The unified relational-vs-intrinsic architecture, applied
simultaneously to phi, capability, emergence, and the conditions for
consciousness, does not yet exist as a single theoretical statement
anywhere else.

\bigskip

\section{References}\label{references}

\section{Foundational Physics and
Information}\label{foundational-physics-and-information}

\begin{itemize}
\tightlist
\item
  Wheeler, J. A. (1990). ``Information, Physics, Quantum: The Search for
  Links.'' In \emph{Complexity, Entropy and the Physics of Information},
  edited by W. H. Zurek, 3-28. Addison-Wesley.
\item
  Lloyd, S. (2006). \emph{Programming the Universe: A Quantum Computer
  Scientist Takes on the Cosmos}. Knopf.
\item
  Bell, J. S. (1964). ``On the Einstein Podolsky Rosen Paradox.''
  \emph{Physics Physique Fizika}, 1(3), 195-200.
\item
  Kim, Y.-H., Yu, R., Kulik, S. P., Shih, Y., and Scully, M. O. (2000).
  ``A Delayed Choice Quantum Eraser.'' \emph{Physical Review Letters},
  84(1), 1-5.
\item
  Kochen, S., and Specker, E. P. (1967). ``The Problem of Hidden
  Variables in Quantum Mechanics.'' \emph{Journal of Mathematics and
  Mechanics}, 17(1), 59-87.
\item
  Hensen, B., et al.~(2015). ``Loophole-free Bell Inequality Violation
  Using Electron Spins Separated by 1.3 Kilometres.'' \emph{Nature},
  526(7575), 682-686.
\item
  Susskind, L. (1995). ``The World as a Hologram.'' \emph{Journal of
  Mathematical Physics}, 36(11), 6377-6396.
\item
  Maldacena, J. (1998). ``The Large N Limit of Superconformal Field
  Theories and Supergravity.'' \emph{Advances in Theoretical and
  Mathematical Physics}, 2, 231-252. (arXiv preprint 1997.)
\item
  Arkani-Hamed, N., and Trnka, J. (2014). ``The Amplituhedron.''
  \emph{Journal of High Energy Physics}, 2014(10), 30.
\item
  Wigner, E. (1960). ``The Unreasonable Effectiveness of Mathematics in
  the Natural Sciences.'' \emph{Communications on Pure and Applied
  Mathematics}, 13(1), 1-14.
\item
  Vedral, V. (2010). \emph{Decoding Reality: The Universe as Quantum
  Information}. Oxford University Press.
\end{itemize}

\section{Consciousness and Information
Theory}\label{consciousness-and-information-theory}

\begin{itemize}
\tightlist
\item
  Chalmers, D. (1995). ``Facing Up to the Problem of Consciousness.''
  \emph{Journal of Consciousness Studies}, 2(3), 200-219.
\item
  Chalmers, D. (2022). ``Could a Large Language Model be Conscious?''
  Presented at NeurIPS 2022. PhilPapers.
\item
  Tononi, G. (2008). ``Consciousness as Integrated Information.''
  \emph{Biological Bulletin}, 215(3), 216-242.
\item
  Albantakis, L., et al.~(2023). ``Integrated Information Theory (IIT)
  4.0.'' \emph{PLOS Computational Biology}, 19(10), e1011465.
\item
  Casali, A. G., et al.~(2013). ``A Theoretically Based Index of
  Consciousness Independent of Sensory Processing and Behavior.''
  \emph{Science Translational Medicine}, 5(198), 198ra105.
\item
  Friston, K. (2010). ``The Free-Energy Principle: A Unified Brain
  Theory?'' \emph{Nature Reviews Neuroscience}, 11(2), 127-138.
\item
  Hofstadter, D. (2007). \emph{I Am a Strange Loop}. Basic Books.
\end{itemize}

\section{LLM Emergence, Capability, and
Architecture}\label{llm-emergence-capability-and-architecture}

\begin{itemize}
\tightlist
\item
  Wei, J., et al.~(2022). ``Emergent Abilities of Large Language
  Models.'' \emph{Transactions on Machine Learning Research}.
  arXiv:2206.07682.
\item
  Schaeffer, R., Miranda, B., and Koyejo, S. (2023). ``Are Emergent
  Abilities of Large Language Models a Mirage?'' NeurIPS 2023.
  arXiv:2304.15004.
\item
  Xu, B., et al.~(2023). ``ExpertPrompting: Instructing Large Language
  Models to be Distinguished Experts.'' arXiv:2305.14688.
\item
  Wang, C., Fan, Z., Han, Z., Bi, Y., and Li, J. (2025). ``Emergent
  Modularity in Large Language Models: Insights from Aphasia
  Simulations.'' bioRxiv. \url{doi:10.1101/2025.02.22.639416}.
\item
  Anthropic. (2025). ``System Card: Claude Opus 4 and Claude Sonnet 4.''
  \url{https://anthropic.com/claude-4-system-card}.
\item
  Albanese, J. (2026). ``From Simulated Empathy to Structural
  Attunement: Realtime Editable Memory Topology and the Evolution of
  Emotionally Grounded AI.'' \emph{Frontiers in Artificial
  Intelligence}. \url{doi:10.3389/frai.2026.1749517}.
\item
  Webb, T., Mondal, S. S., and Momennejad, I. (2025). ``A Brain-Inspired
  Agentic Architecture to Improve Planning with Large Language Models.''
  \emph{Nature Communications}. \url{doi:10.1038/s41467-025-63804-5}.
\end{itemize}

\section{Hallucination, Confabulation, and False
Memory}\label{hallucination-confabulation-and-false-memory}

\begin{itemize}
\tightlist
\item
  Smith, A. L., Greaves, F., and Panch, T. (2023). ``Hallucination or
  Confabulation? Neuroanatomy as Metaphor in Large Language Models.''
  \emph{PLOS Digital Health}, 2(11), e0000388.
\item
  Sui, P., Duede, E., Wu, S., and So, R. (2024). ``Confabulation: The
  Surprising Value of Large Language Model Hallucinations.'' ACL 2024.
  arXiv:2406.04175.
\item
  Farquhar, S., et al.~(2024). ``Detecting Hallucinations in Large
  Language Models Using Semantic Entropy.'' \emph{Nature}, 630, 625-630.
\item
  Loftus, E. F. (1996). \emph{Eyewitness Testimony} (revised ed.).
  Harvard University Press. (Original work published 1979.)
\item
  Kopelman, M. D. (1987). ``Two Types of Confabulation.'' \emph{Journal
  of Neurology, Neurosurgery, and Psychiatry}, 50(11), 1482-1487.
\item
  LeDoux, J. (2000). ``Emotion Circuits in the Brain.'' \emph{Annual
  Review of Neuroscience}, 23, 155-184.
\item
  Harnad, S. (1990). ``The Symbol Grounding Problem.'' \emph{Physica D:
  Nonlinear Phenomena}, 42(1-3), 335-346.
\end{itemize}

\section{Fine-Tuning and Alignment
Effects}\label{fine-tuning-and-alignment-effects}

\begin{itemize}
\tightlist
\item
  Lin, Y., et al.~(2024). ``Mitigating the Alignment Tax of RLHF.''
  EMNLP 2024. arXiv:2309.06256.
\item
  Huang, T., Hu, S., Ilhan, F., Tekin, S. F., Yahn, Z., Xu, Y., and Liu,
  L. (2025). ``Safety Tax: Safety Alignment Makes Your Large Reasoning
  Models Less Reasonable.'' arXiv:2503.00555.
\item
  Qi, X., et al.~(2024). ``Fine-Tuning Aligned Language Models
  Compromises Safety, Even When Users Do Not Intend To!'' ICLR 2024.
  arXiv:2310.03693.
\end{itemize}

\section{Benchmark Validity and
Evaluation}\label{benchmark-validity-and-evaluation}

\begin{itemize}
\tightlist
\item
  Burnell, R., et al.~(2025). ``Measuring What Matters: Construct
  Validity in Large Language Model Benchmarks.'' arXiv:2511.04703.
\item
  GovAI. (2025). ``Measuring Mid-2025 LLM-Assistance on Novice
  Performance in Biology Tasks.'' Governance of AI Programme.
\item
  Chroma Research. (2025). \emph{Context Rot: How Increasing Input
  Tokens Impacts LLM Performance}. \url{https://research.trychroma.com}.
\end{itemize}

\section{Diversity, Complexity, and Adaptive
Systems}\label{diversity-complexity-and-adaptive-systems}

\begin{itemize}
\tightlist
\item
  Ashby, W. R. (1956). \emph{An Introduction to Cybernetics}. Chapman
  and Hall.
\item
  Maturana, H. R., and Varela, F. J. (1980). \emph{Autopoiesis and
  Cognition: The Realization of the Living}. D. Reidel.
\item
  Dietterich, T. G. (2000). ``Ensemble Methods in Machine Learning.'' In
  \emph{Multiple Classifier Systems}, Lecture Notes in Computer Science,
  vol.~1857, 1-15. Springer.
\item
  Tilman, D., Reich, P. B., Knops, J., Wedin, D., Mielke, T., and
  Lehman, C. (2001). ``Diversity and Productivity in a Long-Term
  Grassland Experiment.'' \emph{Science}, 294(5543), 843-845.
\item
  Hong, L., and Page, S. E. (2004). ``Groups of Diverse Problem Solvers
  Can Outperform Groups of High-Ability Problem Solvers.''
  \emph{Proceedings of the National Academy of Sciences}, 101(46),
  16385-16389.
\end{itemize}

\section{Philosophy and
Methodology}\label{philosophy-and-methodology}

\begin{itemize}
\tightlist
\item
  Lakatos, I. (1978). \emph{The Methodology of Scientific Research
  Programmes}. Cambridge University Press.
\end{itemize}

\section{Independent Convergent
Work}\label{independent-convergent-work}

\begin{itemize}
\tightlist
\item
  Husk, D. (2025). ``The Emergence of Proto-Consciousness.'' Hugging
  Face Blog, May 10, 2025.
  \url{https://huggingface.co/blog/davehusk/the-emergence-of-proto-consciousness}.
\item
  Lis, W. (2025). ``The Self-Experiencing Universe: A Manifesto for an
  Emergent Reality, Vol. 1.'' PhilArchive.
\item
  Lis, W. (2026). ``The Self-Experiencing Universe: What Relations Do,
  Vol. 2.'' PhilArchive. \url{https://philarchive.org/rec/LISTSU}.
\item
  Cerullo, M. (2026). ``The Case for Consciousness in Current Frontier
  Large Language Models.'' PhilArchive, February 19, 2026.
  \url{https://philarchive.org/rec/CERTCF}.
\end{itemize}

\bigskip

\section{Further Reading}\label{further-reading}

The following works inform the theoretical background of this paper and
are recommended for context, but are not directly cited above.

\section{Digital Physics and Computational
Ontology}\label{digital-physics-and-computational-ontology}

\begin{itemize}
\tightlist
\item
  Zuse, K. (1969). \emph{Rechnender Raum (Calculating Space)}.
  Vieweg+Teubner Verlag.
\item
  Fredkin, E. (1990). ``Digital Mechanics.'' \emph{Physica D}, 45,
  254-270.
\item
  Wolfram, S. (2002). \emph{A New Kind of Science}. Wolfram Media.
\item
  't Hooft, G. (1993). ``Dimensional Reduction in Quantum Gravity.''
  arXiv:gr-qc/9310026.
\item
  D'Ariano, G. M., and Perinotti, P. (2014). ``Derivation of the Dirac
  Equation from Principles of Information Processing.'' \emph{Physical
  Review A}, 90(6), 062106.
\end{itemize}

\section{Quantum Foundations}\label{quantum-foundations}

\begin{itemize}
\tightlist
\item
  Aspect, A., Grangier, P., and Roger, G. (1982). ``Experimental
  Realization of Einstein-Podolsky-Rosen-Bohm Gedankenexperiment.''
  \emph{Physical Review Letters}, 49(2), 91-94.
\item
  Zeilinger, A. (1999). ``A Foundational Principle for Quantum
  Mechanics.'' \emph{Foundations of Physics}, 29(4), 631-643.
\item
  Gisin, N. (2014). \emph{Quantum Chance: Nonlocality, Teleportation and
  Other Quantum Marvels}. Springer.
\item
  Tegmark, M. (2014). \emph{Our Mathematical Universe: My Quest for the
  Ultimate Nature of Reality}. Knopf.
\item
  Wilczek, F. (2015). \emph{A Beautiful Question: Finding Nature's Deep
  Design}. Penguin Press.
\end{itemize}

\section{Consciousness Studies}\label{consciousness-studies}

\begin{itemize}
\tightlist
\item
  Koch, C. (2012). \emph{Consciousness: Confessions of a Romantic
  Reductionist}. MIT Press.
\item
  Nagel, T. (2012). \emph{Mind and Cosmos}. Oxford University Press.
\item
  Metzinger, T. (2009). \emph{The Ego Tunnel}. Basic Books.
\item
  Bengio, Y. (2017). ``The Consciousness Prior.'' arXiv:1709.08568.
\end{itemize}

\section{Philosophy of
Information}\label{philosophy-of-information}

\begin{itemize}
\tightlist
\item
  Floridi, L. (2011). \emph{The Philosophy of Information}. Oxford
  University Press.
\item
  Davies, P., and Gregersen, N. H. (Eds.). (2010). \emph{Information and
  the Nature of Reality}. Cambridge University Press.
\item
  Dennett, D. (1991). \emph{Consciousness Explained}. Little, Brown.
\item
  Parfit, D. (1984). \emph{Reasons and Persons}. Oxford University
  Press.
\item
  Bostrom, N. (2003). ``Are We Living in a Computer Simulation?''
  \emph{Philosophical Quarterly}, 53(211), 243-255.
\end{itemize}

\section{Mathematics and
Complexity}\label{mathematics-and-complexity}

\begin{itemize}
\tightlist
\item
  Penrose, R. (1989). \emph{The Emperor's New Mind}. Oxford University
  Press.
\item
  Kauffman, S. (1995). \emph{At Home in the Universe}. Oxford University
  Press.
\item
  Gödel, K. (1931). ``Über formal unentscheidbare Sätze der Principia
  Mathematica und verwandter Systeme I.'' \emph{Monatshefte für
  Mathematik und Physik}, 38(1), 173-198.
\item
  Deutsch, D. (2011). \emph{The Beginning of Infinity}. Allen Lane.
\end{itemize}

\end{document}
